Puji syukur ke hadirat Allah SWT atas segala rahmat dan karunia-Nya sehingga buku ``Handout ONMIPA-PT Bidang Matematika: Aljabar Linear'' ini dapat disusun dengan baik oleh penulis dan rekan sejawatnya (GPT Plus).

Buku ini disusun sebagai salah satu bahan pembelajaran dan persiapan dalam menghadapi Olimpiade Nasional Matematika dan Ilmu Pengetahuan Alam Perguruan Tinggi (ONMIPA-PT) Bidang Matematika, yang lomba tersebut pernah penulis ikuti selama 3 kali berturut-turut pada tahun 2024, 2025, dan 2026. Walaupun hanya mentok jadi finalis, setidaknya banyak pengalaman yang dapat penulis ambil dan rasakan hingga mendapatkan niat untuk menyusun buku ini.

Alasan utama penyusunan buku ini adalah untuk merangkum berbagai konsep, definisi, teorema, dan hasil-hasil penting dalam Aljabar Linear yang relevan dengan persiapan ONMIPA-PT. Selain itu, pada dasarnya penulis memang menyukai Aljabar Linear hingga ke tingkat seperti Analisis Fungsional (mata kuliah yang sangat seru sekali hehe). Sangat \textit{recommended} untuk diambil karena campuran antara Aljabar Linear dan Analisis Real yang mungkin saja akan membantu fundamental pembaca untuk menghadapi ONMIPA.

Sembari menunggu \textit{reset token}, penulis menyadari bahwa buku ini masih memiliki berbagai keterbatasan dan jauh dari kesempurnaan, baik dari segi kelengkapan materi, penyajian, maupun penulisan. Oleh karena itu, kritik dan saran yang membangun sangat diharapkan sebagai bahan evaluasi dan perbaikan pada penyusunan berikutnya.

Akhir kata, penulis berharap buku ini dapat memberikan manfaat bagi para pembaca serta menjadi salah satu sarana pendukung dalam proses pembelajaran dan persiapan menghadapi ONMIPA-PT. Semoga buku ini dapat membantu menumbuhkan pemahaman yang lebih mendalam sekaligus meningkatkan minat dan semangat dalam mempelajari matematika.
